\documentclass[11pt]{article}

\usepackage[a4paper,margin=1in]{geometry}
\usepackage{cmap}
\usepackage[T1]{fontenc}
\usepackage[utf8]{inputenc}
\input{glyphtounicode}
\pdfgentounicode=1
\usepackage{lmodern}
\usepackage[ngerman]{babel}
\usepackage{amsmath,amssymb,amsthm}
\usepackage{booktabs}
\usepackage{enumitem}
\usepackage[unicode,pdfencoding=auto,psdextra]{hyperref}

\hypersetup{
  pdftitle={Kein positiver Spielraum aus Büschelargumenten: Xiaos Lemma, eine Spielraumidentität in Rang r und die Nef-Barriere in Höhenargumenten nach Parshin und Moret-Bailly},
  pdfauthor={Lukas Geiger},
  pdfsubject={Büschelargumente für obere Schranken arithmetischer Selbstschnittzahlen},
  pdfkeywords={Parshin-Moret-Bailly, arithmetische Flächen, Xiaos Lemma, Harder-Narasimhan-Filtrierung, Bogomolov-Ungleichung, Kodaira-Faserung},
  bookmarksnumbered=true,
  bookmarksopen=true,
  pdflang={de-DE},
  colorlinks=true,
  linkcolor=black,
  urlcolor=blue,
  citecolor=black
}

\clubpenalty=10000
\widowpenalty=10000
\displaywidowpenalty=10000
\emergencystretch=1.5em

\newtheorem{theorem}{Satz}[section]
\newtheorem{proposition}[theorem]{Proposition}
\newtheorem{lemma}[theorem]{Lemma}
\newtheorem{corollary}[theorem]{Korollar}
\newtheorem*{theoremA}{Satz A}
\newtheorem*{theoremB}{Satz B}
\newtheorem*{theoremC}{Satz C}
\theoremstyle{definition}
\newtheorem{definition}[theorem]{Definition}
\newtheorem{problem}[theorem]{Offenes Problem}
\theoremstyle{remark}
\newtheorem{remark}[theorem]{Bemerkung}

\DeclareMathOperator{\rk}{rk}
\DeclareMathOperator{\Pic}{Pic}
\DeclareMathOperator{\ch}{ch}
\newcommand{\OO}{\mathcal{O}}
\newcommand{\cI}{\mathcal{I}}
\newcommand{\cL}{\mathcal{L}}
\newcommand{\cF}{\mathcal{F}}
\newcommand{\ob}{\overline{\omega}}
\newcommand{\Lb}{\overline{L}}
\newcommand{\degh}{\widehat{\deg}}
\newcommand{\bpf}{\textup{(BPF)}}

\title{Kein positiver Spielraum aus Büschelargumenten:\\ Xiaos Lemma, eine Spielraumidentität in Rang~$r$\\ und die Nef-Barriere in Höhenargumenten nach Parshin und Moret-Bailly}
\author{Lukas Geiger\\
Unabhängiger Forscher, Bernau im Schwarzwald, Deutschland\\
ORCID: \href{https://orcid.org/0009-0005-7296-1534}{0009-0005-7296-1534}}
\date{Entwurf v0.1, 7.~November 2026}

\begin{document}

\maketitle

\begin{abstract}
Höhenargumente nach Parshin und Moret-Bailly zielen auf eine obere Schranke für
die arithmetische Selbstschnittzahl $\ob_P^{\,2}$ der relativ dualisierenden Garbe
auf den Fasern einer Faserung $\phi\colon X\to Y$ über einer Kurve. Ein neuer
Ansatz dieser Art (Soul\'e, arXiv:2611.05217, Version~1) gewinnt die
entscheidende Abschätzung aus einem Büschel, nämlich aus zwei Schnitten eines
Hilfsgeradenbündels $L$ ohne gemeinsame Nullstelle, dem Satz von Minkowski über
sukzessive Minima und einem Übergang von klassischen zu arithmetischen Graden. Der
führende Koeffizient der resultierenden Ungleichung ist der geometrische
Spielraum $M=(-L^2+m\deg\phi_*L)/\omega^2$, wobei $m$ den Fasergrad von $L$
bezeichnet; eine obere Schranke für $\ob_P^{\,2}$ erfordert $M>0$. Wir zeigen,
dass dieser Fall nie eintritt. In Rang zwei liefert die Koszul-Auflösung der
Auswertungsabbildung die exakte Identität
$-L^2+m\deg\phi_*L=-\ell(Z)-2m\deg V\le 0$ (Xiaos Lemma für Büschel), wobei $Z$
und $V$ den endlichen bzw.\ den vertikalen Anteil des Basisortes bezeichnen. In
Rang $r$ beweisen wir eine Spielraumidentität mittels der
Bogomolov-Diskriminante des Syzygienkerns sowie ein Nef-Lemma im Sinne Xiaos;
dieses liefert $M_r\le 0$, sobald $\phi_*L$ semistabil ist, und strikte oder
partielle Negativität im Harder--Narasimhan-instabilen Fall. Die arithmetische
Seite des instabilen Falls bleibt für $r\ge 3$ offen. Als Anwendung zeigen wir,
dass Proposition~9\,vi) in arXiv:2611.05217v1 falsch ist: Für die dort verwendete
Kodaira-Faserung gilt $k=4$, über jedem exzeptionellen Punkt sind höchstens zwei
der vier Quadratwurzeln effektiv, und es folgt $\deg\phi_*L=-7k+4k_{\mathrm{eff}}\le
-20<-3k$. Der fehlerhafte Beweisschritt verwendet die Umkehrung der in
Proposition~5\,ii) jenes Preprints bewiesenen Implikation. Abschließend gehen wir
kurz auf Kreuzterme, einzelne Schnitte und Jets längs Kurven sowie auf Quellbündel
ein, die nicht von der Basis zurückgezogen sind. Über die $abc$-Vermutung wird
nichts ausgesagt.
\end{abstract}

\section{Einleitung}\label{sec:intro}

Parshin~\cite{Pa88} schlug ein arithmetisches Analogon der
Bogomolov--Miyaoka--Yau-Ungleichung für komplexe Flächen vor (zu den
Ungleichungen von Bogomolov und Miyaoka siehe~\cite{Bo78,Mi77}): eine obere
Schranke für die arithmetische Selbstschnittzahl der relativ dualisierenden Garbe
einer arithmetischen Fläche in Abhängigkeit von der Diskriminante der Basis.
Parshin und Moret-Bailly~\cite{Pa88,MB90} zeigten, dass eine solche Schranke,
angewandt auf die Fasern einer geeigneten Faserung, weitreichende diophantische
Konsequenzen hat, darunter effektive Aussagen vom Mordell-Typ und Formen der
$abc$-Vermutung (siehe auch~\cite{So05}). Argumente dieser Art nennen wir
\emph{Höhenargumente nach Parshin und Moret-Bailly (PMB-Argumente)}.

In einem neuen Preprint~\cite{So26} (arXiv:2611.05217, Version~1, eingereicht am
3.~November 2026) setzt Soul\'e diese Strategie für eine Kodaira-Faserung
$\phi\colon X\to Y$ mit Fasern vom Geschlecht~$6$ über einer étalen Überlagerung
einer Kurve vom Geschlecht~$2$ um~\cite{Ko67,MD85}. Die obere Schranke für
$\ob_P^{\,2}$ auf der Faser $X_P$ beruht auf einem \emph{Büschel}: einem
Hilfsgeradenbündel $L$ auf $X$, zwei Schnitten von $L_P$ ohne gemeinsame
Nullstelle~\cite[Prop.~3]{So26}, dem Satz von Minkowski über sukzessive Minima
und einem Übergang von klassischen zu arithmetischen Graden~\cite[Cor.~1]{So26}.
Der führende Koeffizient der resultierenden Ungleichung ist eine rein geometrische
Größe, die wir den \emph{Spielraum} nennen:
\begin{equation}\label{eq:margin-intro}
M(L)=\frac{-L^2+m\deg\phi_*L}{\omega^2},
\end{equation}
wobei $m$ den Grad von $L$ auf einer Faser bezeichnet und $\omega=\omega_{X/Y}$
gesetzt ist. Das Argument beschränkt $\ob_P^{\,2}$ genau dann nach oben, wenn
$M(L)>0$ gilt (Abschnitt~\ref{sec:setting}). In~\cite{So26} stützt sich die
Positivität von $M(L)$ auf die untere Schranke $\deg\phi_*L\ge -3k$
aus~\cite[Prop.~9\,vi)]{So26}.\footnote{Wir zitieren~\cite{So26} nach der
Nummerierung der Druckfassung von Version~1. In der \TeX-Quelle von Version~1
tragen die als Propositionen~3, 4, 5, 8 und~9 gedruckten Aussagen die Labels
\texttt{prop12}, \texttt{prop3}, \texttt{prop13}, \texttt{prop5} und
\texttt{prop6}, während Lemma~3 das Label \texttt{lm15} trägt. Gleichungsnummern
wie (24) und (28)--(36) sind in der Quelle manuell vergeben und stimmen mit der
Druckfassung überein.}

In dieser Note zeigen wir, dass der Spielraum eines Büschels nie positiv ist, und
ziehen daraus die Folgerung für~\cite{So26}.

\begin{theoremA}[Satz~\ref{thm:rank2}]
Es sei $\phi\colon X\to Y$ eine Faserung einer glatten projektiven Fläche über
einer glatten projektiven Kurve mit integren Fasern und $L$ ein Geradenbündel mit
$h^0(L_\eta)=2$, dessen Linearsystem auf der generischen Faser basispunktfrei
ist. Bezeichnet $L(-V)\otimes\cI_Z$ das Bild von $\phi^*\phi_*L\to L$, so gilt
\[
-L^2+m\deg\phi_*L=-\ell(Z)-2m\deg V\le 0 .
\]
Folglich ist $M(L)\le 0$, mit Gleichheit nur bei surjektiver Auswertungsabbildung.
\end{theoremA}

Dies ist Xiaos Lemma für Büschel~\cite{Xi87} in exakter Form; eine
Stabilitätsvoraussetzung ist hierfür nicht erforderlich.

\begin{theoremB}[Satz~\ref{thm:rankr}, Korollar~\ref{cor:ss}, Proposition~\ref{prop:unstable}]
Für $r=\rk\phi_*L\ge 2$ lautet der maßgebliche Spielraum
$M_r(L)=(-L^2+\tfrac{2m}{r}\deg\phi_*L)/\omega^2$, und es gilt
\[
M_r(L)=-\frac{\Delta(K)+2rm\deg V+2(r-1)\ell(Z)}{r\,\omega^2},
\]
wobei $\Delta(K)$ die Bogomolov-Diskriminante des Syzygienkerns $K$ bezeichnet.
Ist $\phi_*L$ semistabil, so gilt $M_r(L)\le 0$. Ist $\phi_*L$ instabil, so gilt
$M_r(L)<0$, sofern die maximal destabilisierende Untergarbe keinen horizontalen
Basisdivisor besitzt; allgemein ist jeder wie in
Proposition~\ref{prop:unstable}\,(b) aus zwei Harder--Narasimhan-Steigungen
gebildete Koeffizient nichtpositiv.
\end{theoremB}

Der semistabile Fall stützt sich auf ein Nef-Lemma im Sinne der Methode von Xiao:
Die der Harder--Narasimhan-Filtrierung von $\phi_*L$ zugeordneten
$\mathbb{Q}$-Divisoren $L-B_j-V_j-\mu_jF$ sind nef (Lemma~\ref{lem:nef}). Das
arithmetische Gegenstück des instabilen Falls bleibt für $r\ge 3$ offen (Offenes
Problem~\ref{prob:anomaly}).

\begin{theoremC}[Sätze~\ref{thm:soule-count} und~\ref{thm:soule-structural}]
Für die Kodaira-Faserung aus~\cite{So26} gilt $k=4$; über jedem exzeptionellen
Punkt sind höchstens zwei der vier Quadratwurzeln der hyperelliptischen Klasse
effektiv, und es gilt
\[
\deg\phi_*L=-7k+4k_{\mathrm{eff}}\le -20<-12=-3k ,
\]
wobei $4k_{\mathrm{eff}}\le 8$ die Anzahl der Punkte von $Y$ bezeichnet, über
denen die jeweilige Quadratwurzel effektiv ist. Somit ist~\cite[Prop.~9\,vi)]{So26}
falsch, und der führende Koeffizient in~\cite[(36)]{So26} beträgt höchstens
$-5/6$. Ein zweiter Zugang liefert dasselbe Ergebnis: Das generische Linearsystem
ist für jede Wahl der Komponente $Y$ und der Quadratwurzel $\cL$ basispunktfrei,
sodass Satz~A die Schranke $\deg\phi_*L\le -10k/3<-3k$ ergibt;
Basispunktfreiheit ist genau jene Eigenschaft, die in~\cite[(28)]{So26} benutzt
wird, um~\cite[Prop.~3]{So26} anzuwenden.
\end{theoremC}

Der fehlerhafte Schritt im Beweis von~\cite[Prop.~9]{So26} liegt in
Ungleichung~(24) jener Arbeit: Dort wird die Implikation „$P_1$ exzeptionell
$\Rightarrow
h^0(\lambda)=1$“ verwendet, während~\cite[Prop.~5\,ii)]{So26} die Umkehrung
beweist (Bemerkung~\ref{rem:gap}). Abschnitt~\ref{sec:scope} legt kurz dar,
weshalb dasselbe Vorzeichen auch bei Kreuztermen sowie bei einzelnen Schnitten und
Jets längs Kurven auftritt und weshalb ein positiver Beitrag nur von einem Bündel
herrühren kann, das nicht von der Basis zurückgezogen ist.

\paragraph{Geltungsgrenze.}
Diese Note trifft keine Aussage über die $abc$-Vermutung, über Parshins Vermutung
oder über die Gültigkeit der Sätze~1, A oder~B aus~\cite{So26} als Aussagen. Sie
zeigt, dass der Büschelschritt eines PMB-Arguments unter den unten genannten
Voraussetzungen keinen positiven Spielraum liefern kann und dass der in Version~1
gegebene Beweis von~\cite[Theorem~1]{So26} bei Proposition~9\,vi) unvollständig
ist. Alle Verweise auf~\cite{So26} betreffen Version~1; eine spätere Version kann
abweichen. Die von uns herangezogenen Schnittzahlen
aus~\cite[Prop.~8 und Prop.~9\,ii), iii)]{So26} stimmen mit unserer eigenen
Rechnung überein.

\section{Rahmen und Spielraum eines Büschels}\label{sec:setting}

\subsection{Konventionen}
In den Abschnitten~\ref{sec:setting}--\ref{sec:nef} arbeiten wir über einem
algebraisch abgeschlossenen Körper der Charakteristik~$0$. Es sei $\phi\colon X\to
Y$ ein surjektiver Morphismus einer glatten projektiven Fläche auf eine glatte
projektive Kurve, dessen Fasern sämtlich integer sind. Ferner bezeichne $F$ die
Klasse einer Faser, $\eta$ den generischen Punkt von~$Y$ und
$\omega=\omega_{X/Y}$, wobei $\omega^2>0$ gelte. Glatte Faserungen, insbesondere
Kodaira-Faserungen, erfüllen diese Voraussetzungen. Ein Divisor heißt
\emph{vertikal}, wenn seine Komponenten in Fasern liegen, und \emph{horizontal},
wenn keine Komponente in einer Faser liegt. Da die Fasern integer sind, hat jeder
effektive vertikale Divisor die Gestalt $V=\sum_y n_yX_y$ mit $n_y\ge 0$; wir
setzen $\deg V=\sum_y n_y$, sodass $V\equiv(\deg V)F$, $V^2=0$ und $V\cdot F=0$
gelten.

Für ein Geradenbündel $L$ auf $X$ sei $m=L\cdot F\ge 1$ und $E=\phi_*L$. Als
torsionsfreie Garbe auf der glatten Kurve $Y$ ist $E$ lokal frei vom Rang
$r=h^0(X_\eta,L_\eta)$; wir schreiben $\mu(E)=\deg E/r$. Es gilt $L\cdot V=m\deg
V$. Häufig setzen wir voraus:
\begin{itemize}[leftmargin=3.5em]
\item[\bpf] Das Linearsystem $|L_\eta|$ auf der generischen Faser hat keinen
Basispunkt.
\end{itemize}

\begin{lemma}\label{lem:image}
Sei $0\ne E'\subseteq E$ eine Untergarbe. Das Bild der Komposition
$\phi^*E'\to\phi^*E\to L$ hat die Gestalt $L(-B'-V')\otimes\cI_{Z'}$ mit einem
horizontalen Divisor $B'\ge 0$, einem vertikalen Divisor $V'\ge 0$ und einem
Unterschema $Z'$ endlicher Länge. Die Einschränkung von $L(-B')$ auf $X_\eta$
besitzt einen von null verschiedenen Schnitt. Unter~\bpf{} gilt $B'=0$ für
$E'=E$.
\end{lemma}

\begin{proof}
Auf einer glatten Fläche hat jede von null verschiedene Untergarbe eines
Geradenbündels $L$ die Gestalt $L(-D)\otimes\cI_{Z'}$ mit einem effektiven
Divisor $D\ge 0$ und einem Unterschema $Z'$ endlicher Länge; wir zerlegen
$D=B'+V'$ in seinen horizontalen und seinen vertikalen Anteil. Über $\eta$ fällt
die Abbildung mit der Auswertung der Schnitte in
$E'_\eta\subseteq E_\eta=H^0(X_\eta,L_\eta)$ zusammen; ihr Bild ist die von diesen
Schnitten erzeugte Untergarbe mit divisoriellem Basisort $B'|_{X_\eta}$, und jeder
von null verschiedene Schnitt aus $E'_\eta$ liefert einen Schnitt von
$L(-B')|_{X_\eta}$. Für $E'=E$ ist dieser Basisort unter~\bpf{} leer.
\end{proof}

\begin{definition}\label{def:margin}
Es sei $r_1=-L^2/\omega^2$ und $q=-\deg E/\omega^2$. Der \emph{Spielraum} von
$L$ ist definiert durch
\[
M_r(L)=r_1-\frac{2m}{r}\,q=\frac{-L^2+\frac{2m}{r}\deg E}{\omega^2}.
\]
Für $r=2$ schreiben wir $M(L)=M_2(L)=r_1-mq$, in Übereinstimmung
mit~\eqref{eq:margin-intro}.
\end{definition}

\begin{lemma}\label{lem:square}
Es gilt $M_r(L)=-(L-\mu(E)F)^2/\omega^2$ sowie $M_r(L+dF)=M_r(L)$ für jede ganze
Zahl $d$.
\end{lemma}

\begin{proof}
Es gilt $(L-\mu F)^2=L^2-2m\mu=L^2-\frac{2m}{r}\deg E$. Ein Twisten mit $dF$
ändert $L^2$ um $2md$ und $\deg E$ nach der Projektionsformel um $rd$.
\end{proof}

\subsection{Die Büschelkette}\label{subsec:chain}
Es sei nun $\phi\colon X\to Y$ über einem Zahlkörper~$K$ definiert und habe
geometrisch integre Fasern; die obigen Konventionen wenden wir auf
$X_{\overline K}\to Y_{\overline K}$ an, wodurch sich Schnittzahlen und Grade
nicht ändern. Für eine endliche Erweiterung $F/K$ vom Grad $d$ und einen Punkt
$P\in Y(F)$ bezeichne $X_P$ die Faser; wir versehen $\omega$ und $L$ wie
in~\cite[\S1.3]{So26} mit adelischen Metriken und $H_P=H^0(X_P,L_P)$ mit den
Supremumsnormen. Für $r=2$ verbindet der Beweis von~\cite[Theorem~1]{So26}
in~\cite[\S6.3]{So26} drei Schritte:
\begin{enumerate}[label=(\roman*)]
\item \emph{Büschelungleichung}~\cite[Prop.~3, (28), (29)]{So26}: Haben
$s_1,s_2\in H_P$ keine gemeinsame Nullstelle, so gilt
$-\Lb_P^{\,2}\le m\,(h(s_1)+h(s_2))+O(d)$.
\item \emph{Minkowski}~\cite[(2), (30)]{So26}: Es gibt linear unabhängige
$s_1,s_2\in H_P$ mit $h(s_1)+h(s_2)\le-\degh(H_P)+\delta(F)+O(d)$.
\item \emph{Übertragung}~\cite[(32), Cor.~1, (33), (35)]{So26}: Für $P$
außerhalb einer endlichen Menge und jedes $\varepsilon>0$ gilt
$\degh(H_P)=-q\,\ob_P^{\,2}+O(\varepsilon\,\ob_P^{\,2})+O_\varepsilon(d)$ und
$\Lb_P^{\,2}=-r_1\,\ob_P^{\,2}+O(\varepsilon\,\ob_P^{\,2})+O_\varepsilon(d)$.
\end{enumerate}
Zusammen ergibt sich
\begin{equation}\label{eq:chain}
\bigl(M(L)-O(\varepsilon)\bigr)\,\ob_P^{\,2}\le m\,\delta(F)+O_\varepsilon(d).
\end{equation}
Für allgemeines $r$ erfüllen zwei Schnitte, welche die ersten beiden sukzessiven
Minima von $H_P$ realisieren, die Ungleichung
$\lambda_1+\lambda_2\le\frac{2}{r}(\lambda_1+\dots+\lambda_r)$, und der zweite
Satz von Minkowski beschränkt die Summe bis auf ein Vielfaches von $\delta(F)$
und $O(d)$ durch $-\degh(H_P)$. Haben diese beiden Schnitte keine gemeinsame
Nullstelle, so liefern dieselben Schritte~\eqref{eq:chain} mit $M_r(L)$ anstelle
von $M(L)$ und einer angepassten Konstante vor $\delta(F)$.

\begin{remark}\label{rem:chain}
Ungleichung~\eqref{eq:chain} beschränkt $\ob_P^{\,2}$ genau dann nach oben, wenn
$M_r(L)>0$ gilt; für $M_r(L)\le 0$ ist sie inhaltsleer. In~\cite{So26} gilt $r=2$,
$m=6$, $L^2=-20k$ und $\omega^2=12k$, also $r_1=5/3$; Proposition~9\,vi) liefert
$q\le 1/4$, damit $M(L)\ge 1/6$ und die Konstante $36=6/(1/6)$
aus~\cite[Theorem~1]{So26}. Die $\varepsilon$- und $O(d)$-Terme ändern das
Vorzeichen nicht. Alle nachfolgenden Ergebnisse sind Aussagen über die
klassischen Invarianten $L^2$, $\deg\phi_*L$ und $\omega^2$ und in diesem Sinne
unabhängig von den arithmetischen Einzelheiten der Kette.
\end{remark}

\section{Rang zwei: Xiaos Lemma für Büschel}\label{sec:rank2}

\begin{theorem}\label{thm:rank2}
Unter den Voraussetzungen von Abschnitt~\ref{sec:setting} gelte $r=2$ und~\bpf,
und es sei $L(-V)\otimes\cI_Z$ das Bild der Auswertungsabbildung $\phi^*E\to L$
(Lemma~\ref{lem:image}). Dann gilt
\begin{equation}\label{eq:rank2}
-L^2+m\deg E=-\ell(Z)-2m\deg V\le 0 .
\end{equation}
Insbesondere ist $\deg\phi_*L\le L^2/m$, mit Gleichheit genau dann, wenn
$\phi^*\phi_*L\to L$ surjektiv ist.
\end{theorem}

\begin{proof}
Es sei $A=L(-V)$; die Auswertungsabbildung induziert dann eine Surjektion $\phi^*E\to
A\otimes\cI_Z$, deren Kern wir mit $K$ bezeichnen.

\emph{Schritt 1: $K$ ist ein Geradenbündel.} In jedem Punkt ist das Ideal
$\cI_Z$ ein torsionsfreier Modul über einem zweidimensionalen regulären lokalen
Ring; es besitzt daher mindestens die Tiefe eins und nach dem Satz von
Auslander--Buchsbaum höchstens die projektive Dimension eins. Der Kern einer
Surjektion einer lokal freien Garbe auf eine Garbe der projektiven Dimension
höchstens eins ist lokal frei; somit hat $K$ den Rang eins.

\emph{Schritt 2: Chern-Klassen der äußeren Terme.} Da $\phi^*E$ von der Kurve
$Y$ zurückgezogen ist, gelten $c_1(\phi^*E)\equiv(\deg E)F$ und
$c_2(\phi^*E)=\phi^*c_2(E)=0$. Aus der exakten Sequenz
$0\to\cI_Z\to\OO_X\to\OO_Z\to 0$ und $\ch(\OO_Z)=\ell(Z)[\mathrm{pt}]$ folgen
$c_1(\cI_Z)=0$ und $c_2(\cI_Z)=\ell(Z)$; Tensorieren der Rang-eins-Garbe $\cI_Z$
mit $A$ liefert $c_1(A\otimes\cI_Z)=A$ und $c_2(A\otimes\cI_Z)=\ell(Z)$.

\emph{Schritt 3: Whitney-Formel.} Es gilt $c_1(K)\equiv(\deg E)F-A$ sowie
\[
0=c_2(\phi^*E)=c_1(K)\cdot A+\ell(Z)=m\deg E-A^2+\ell(Z),
\]
wobei $F\cdot A=m$ (wegen $F\cdot V=0$) eingeht.

\emph{Schritt 4.} Wegen $V^2=0$ und $L\cdot V=m\deg V$ gilt
$A^2=(L-V)^2=L^2-2m\deg V$. Einsetzen liefert~\eqref{eq:rank2}. Gleichheit
$\deg E=L^2/m$ tritt genau dann ein, wenn $Z=\emptyset$ und $V=0$ gilt, wenn also
die Auswertungsabbildung surjektiv ist.
\end{proof}

\begin{corollary}\label{cor:rank2}
Unter den Voraussetzungen von Satz~\ref{thm:rank2} gilt
\[
M(L)=-\frac{\ell(Z)+2m\deg V}{\omega^2}=-\frac{(L-\mu(E)F)^2}{\omega^2}\le 0 .
\]
Die Büschelkette~\eqref{eq:chain} ist daher für jedes Geradenbündel vom Rang zwei
mit basispunktfreiem generischem Linearsystem auf jeder Faserung mit integren
Fasern inhaltsleer.
\end{corollary}

\begin{remark}\label{rem:rank2}
(a) Satz~\ref{thm:rank2} kommt ohne Stabilitätsvoraussetzung an $\phi_*L$ aus; er
stellt die exakte Rang-zwei-Fassung von Xiaos Methode~\cite{Xi87} dar. Äquivalent
dazu ist $(L-\mu(E)F)^2=\ell(Z)+2m\deg V\ge 0$.

(b) Für $r=2$ hilft auch keine geschicktere Wahl der beiden Schnitte weiter. Der
zweite Satz von Minkowski legt $\lambda_1+\lambda_2$ bis auf Terme der Ordnung
$\delta(F)+d$ fest; zusammen mit Schritt~(iii) erfüllt daher jedes Paar linear
unabhängiger Schnitte die Abschätzung
$h(s_1)+h(s_2)\ge q\,\ob_P^{\,2}-O(\varepsilon\,\ob_P^{\,2}+\delta(F)+d)$. Der
bestmögliche Koeffizient, den die Kette erreichen kann, ist somit
$r_1-mq=M(L)\le 0$.

(c) Auch ein Twisten mit vertikalen Divisoren ändert daran nichts: Nach
Lemma~\ref{lem:square} gilt $M(L+dF)=M(L)$.
\end{remark}

\section{Rang \texorpdfstring{$r$}{r}: eine Spielraumidentität}\label{sec:rankr}

Für eine lokal freie Garbe $G$ vom Rang $n$ auf $X$ bezeichne
$\Delta(G)=2n\,c_2(G)-(n-1)\,c_1(G)^2=c_1(G)^2-2n\ch_2(G)$ ihre
Bogomolov-Diskriminante; für Geradenbündel gilt $\Delta(G)=0$.

\begin{theorem}\label{thm:rankr}
Unter den Voraussetzungen von Abschnitt~\ref{sec:setting} gelte~\bpf{} und
$r\ge 2$. Es sei $L(-V)\otimes\cI_Z$ das Bild von $\phi^*E\to L$ und $K$ der
zugehörige Kern. Dann ist $K$ lokal frei vom Rang $r-1$, und es gilt
\begin{equation}\label{eq:rankr}
\Delta(K)=r\,L^2-2m\deg E-2rm\deg V-2(r-1)\ell(Z),
\end{equation}
oder äquivalent
\[
M_r(L)=-\frac{\Delta(K)+2rm\deg V+2(r-1)\ell(Z)}{r\,\omega^2}.
\]
\end{theorem}

\begin{proof}
Die lokale Freiheit ergibt sich wie in Schritt~1 des Beweises von
Satz~\ref{thm:rank2}. Sei $A=L(-V)$. Aus $c(\phi^*E)=c(K)\,c(A\otimes\cI_Z)$ und
Schritt~2 folgen $c_1(K)\equiv(\deg E)F-A$ sowie
$0=c_2(K)+c_1(K)\cdot A+\ell(Z)$, also
\[
c_1(K)^2=A^2-2m\deg E,\qquad c_2(K)=A^2-m\deg E-\ell(Z).
\]
Damit ist
$\Delta(K)=2(r-1)c_2(K)-(r-2)c_1(K)^2=r\,A^2-2m\deg E-2(r-1)\ell(Z)$, und Einsetzen
von $A^2=L^2-2m\deg V$ ergibt~\eqref{eq:rankr}. Division durch $-r\,\omega^2$
liefert die Darstellung von $M_r(L)$. Für $r=2$ ist der Kern ein Geradenbündel,
es gilt $\Delta(K)=0$, und~\eqref{eq:rankr} reduziert sich auf~\eqref{eq:rank2}.
\end{proof}

\begin{corollary}\label{cor:bogomolov}
Ist $K$ zusätzlich $\mu$-semistabil bezüglich einer amplen Klasse auf $X$, so
liefert die Bogomolov-Ungleichung $\Delta(K)\ge 0$~\cite{Bo78} (siehe
auch~\cite[\S3.4]{HL10}) die Abschätzung $M_r(L)\le 0$.
\end{corollary}

\begin{remark}\label{rem:rankr}
Die Semistabilität des Syzygienkerns ist nicht automatisch gegeben, und für
$r\ge 3$ ist die Abschätzung $L^2\ge\frac{2m}{r}\deg\phi_*L$ eine
Steigungsungleichung, die in der Literatur unter Stabilitätsvoraussetzungen
hergeleitet wird (vgl.~\cite{Xi87,CH88}). Abschnitt~\ref{sec:nef} ersetzt die
Voraussetzung an $K$ durch die Semistabilität von $\phi_*L$ auf der Basiskurve und
behandelt darüber hinaus den instabilen Fall.
\end{remark}

\section{Das Nef-Lemma und die Harder--Narasimhan-Dichotomie}\label{sec:nef}

Es sei $0=E_0\subsetneq E_1\subsetneq\dots\subsetneq E_n=E$ die
Harder--Narasimhan-Filtrierung von $E=\phi_*L$~\cite{HN75} mit den Steigungen
$\mu_j=\mu(E_j/E_{j-1})$; es gilt also $\mu_1>\dots>\mu_n$, und $\mu_j$ ist die
minimale Steigung von $E_j$. Nach Lemma~\ref{lem:image} hat das Bild von
$\phi^*E_j\to L$ die Gestalt $L(-B_j-V_j)\otimes\cI_{Z_j}$; unter~\bpf{} gilt
$B_n=0$, $V_n=V$ und $Z_n=Z$. Wir setzen
\[
N_j=L-B_j-V_j-\mu_jF\in\Pic(X)\otimes\mathbb{Q}.
\]

\begin{lemma}[Nef-Lemma]\label{lem:nef}
Für jedes $j$ ist der $\mathbb{Q}$-Divisor $N_j$ nef. Die
Voraussetzung~\bpf{} wird hierfür nicht benötigt.
\end{lemma}

\begin{proof}
Sei $C\subset X$ eine irreduzible Kurve. Ist $C$ vertikal, so ist $C$ eine Faser
und es gilt $N_j\cdot C=(L-B_j)\cdot F\ge 0$, da $L(-B_j)|_{X_\eta}$ nach
Lemma~\ref{lem:image} einen von null verschiedenen Schnitt besitzt. Sei nun $C$
horizontal, $\nu\colon\widetilde C\to C$ die Normalisierung und
$f=\phi\circ\nu\colon\widetilde C\to Y$; dies ist ein endlicher Morphismus vom
Grad $d_C=C\cdot F\ge 1$. Die Abbildung
$\phi^*E_j\to L(-B_j-V_j)\otimes\cI_{Z_j}$ ist außerhalb der endlichen Menge
$Z_j$ surjektiv auf $L(-B_j-V_j)$, insbesondere also im generischen Punkt von
$C$. Rückzug auf $\widetilde C$ liefert einen von null verschiedenen
Homomorphismus $f^*E_j\to\nu^*L(-B_j-V_j)$, dessen Bild $Q$ ein in
$\nu^*L(-B_j-V_j)$ enthaltener invertierbarer Quotient von $f^*E_j$ ist. Damit
folgt
\[
(L-B_j-V_j)\cdot C=\deg\nu^*L(-B_j-V_j)\ge\deg Q\ge\mu_{\min}(f^*E_j)=d_C\,\mu_j .
\]
Die letzte Gleichheit beruht darauf, dass in Charakteristik~$0$ der Rückzug der
Harder--Narasimhan-Filtrierung unter einem endlichen Morphismus glatter Kurven
mit der Harder--Narasimhan-Filtrierung des Rückzugs übereinstimmt, wobei sich die
Steigungen mit dem Abbildungsgrad multiplizieren (siehe~\cite[\S3.2]{HL10}). Wegen
$F\cdot C=d_C$ ist dies äquivalent zu $N_j\cdot C\ge 0$.
\end{proof}

\begin{corollary}[semistabiler Fall]\label{cor:ss}
Es gelte~\bpf, und $E=\phi_*L$ sei semistabil. Dann ist $N=L-V-\mu(E)F$ nef, und
es gilt
\[
M_r(L)=-\frac{N^2+2m\deg V}{\omega^2}\le 0 .
\]
Eine Voraussetzung an die Stabilität der Fasern, ihrer Linearsysteme oder des
Syzygienkerns ist nicht erforderlich.
\end{corollary}

\begin{proof}
Hier ist $n=1$, $E_1=E$, $B_1=0$, $V_1=V$ und $\mu_1=\mu(E)$; nach
Lemma~\ref{lem:nef} ist $N=N_1$ daher nef, also gilt $N^2\ge 0$. Wegen
$L-\mu(E)F=N+V$, $N\cdot V=m\deg V$ und $V^2=0$ folgt
$(L-\mu(E)F)^2=N^2+2m\deg V$; die Behauptung ergibt sich nun aus
Lemma~\ref{lem:square}.
\end{proof}

\begin{proposition}[instabiler Fall]\label{prop:unstable}
Es gelte~\bpf, und $E=\phi_*L$ sei instabil, $n\ge 2$.
\begin{enumerate}[label=(\alph*)]
\item Ist $B_1=0$, so gilt
$M_r(L)\le-2m\bigl(\mu_1-\mu(E)+\deg V_1\bigr)/\omega^2<0$.
\item Für alle Indizes $a,b$ mit $B_b=0$ gilt
\[
L^2\ge m(\mu_a+\mu_b)+m\deg V_b ,\qquad\text{d.\,h.}\qquad
\frac{-L^2+m(\mu_a+\mu_b)}{\omega^2}\le-\frac{m\deg V_b}{\omega^2}\le 0 .
\]
Ein solcher Index $b$ existiert stets, da $B_n=0$ gilt.
\end{enumerate}
\end{proposition}

\begin{proof}
(a) Es gilt $L-\mu(E)F=N_1+V_1+(\mu_1-\mu(E))F$. Ausmultiplizieren unter
Verwendung von $N_1\cdot V_1=m\deg V_1$, $N_1\cdot F=m$ (wegen $B_1=0$) und
$V_1^2=V_1\cdot F=F^2=0$ ergibt
$(L-\mu(E)F)^2=N_1^2+2m\deg V_1+2m(\mu_1-\mu(E))$. Da $N_1^2\ge 0$ und für
instabiles $E$ zudem $\mu_1>\mu(E)$ gilt, folgt die Behauptung mit
Lemma~\ref{lem:square}.

(b) Die Divisoren $N_b$ und $N_a$ sind nef, $B_a$ und $V_a$ effektiv, also gilt
$N_b\cdot(N_a+B_a+V_a)\ge 0$. Nun ist $N_a+B_a+V_a=L-\mu_aF$ und
$N_b=L-V_b-\mu_bF$, und das Schnittprodukt lautet $L^2-m\mu_a-m\mu_b-m\deg V_b$.
\end{proof}

\begin{remark}\label{rem:unstable}
Im Fall (a) ist der Spielraum, den die Kette aus Abschnitt~\ref{subsec:chain} mit
dem Satz von Minkowski auf dem vollen Gitter liefert, strikt negativ. Die
Koeffizienten in (b) sind genau jene, welche die Kette lieferte, wenn die Höhen
der beiden Schnitte durch die Steigungen $\mu_a$ und $\mu_b$ bestimmt wären;
beispielsweise ergibt Corollary~1 aus~\cite{So26}, angewandt auf das von $E_1$
herrührende Untergitter, einen Schnitt $s_1$ mit
$h(s_1)\le(-\mu_1/\omega^2+\varepsilon)\,\ob_P^{\,2}+c\,\delta(F)+O_\varepsilon(d)$
für eine Konstante $c$. Jeder Koeffizient, den die Kette auf diesem Weg aus
geometrischen Unterbündeln gewinnen kann, ist somit nichtpositiv. Die Bedingung
$B_b=0$ in (b) gibt den Büschelschritt wieder: Jeder von $E_j$ herrührende Schnitt
verschwindet längs $B_j\cap X_P$; daher besitzen zwei Schnitte aus derselben Stufe
$E_a$ mit $B_a\ne 0$ stets eine gemeinsame Nullstelle auf $X_P$. Haben $B_a$ und
$B_b$ lediglich keine gemeinsame Komponente, so sind ihre Einschränkungen auf
$X_P$ für alle bis auf endlich viele $P$ disjunkt, und dieselbe Nichtpositivität
folgt aus Proposition~\ref{prop:cross} mit $L_1=L_2=L$.
\end{remark}

\begin{problem}[arithmetische Anomalie, $r\ge 3$]\label{prob:anomaly}
Es sei $r\ge 3$ und $\phi_*L$ instabil. Es ist offen, ob sich für Punkte $P$
großer Höhe über Paare von Schnitten ohne gemeinsame Nullstelle eine obere
Schranke für $h(s_1)+h(s_2)$ gewinnen lässt, die kleiner ausfällt als der durch
die geometrischen Steigungen $\mu_a$ nahegelegte Wert. Um dies auszuschließen,
müsste man das arithmetische Harder--Narasimhan-Polygon des adelischen Gitters
$H_P$ mit dem um den Faktor $\ob_P^{\,2}/\omega^2$ skalierten geometrischen
Harder--Narasimhan-Polygon von $\phi_*L$ vergleichen. Wir lassen diese Frage
offen. Für $r=2$ tritt sie nicht auf (Bemerkung~\ref{rem:rank2}\,(b)), und für
semistabiles $\phi_*L$ wird sie durch Korollar~\ref{cor:ss} beantwortet.
\end{problem}

\section{Anwendung auf Soul\'e, arXiv:2611.05217v1}\label{sec:soule}

\subsection{Die Konstruktion}
Wir erinnern zunächst an die auf~\cite{Ko67,MD85} zurückgehende Kodaira-Faserung
aus~\cite[\S\S3, 5]{So26}. Es sei $Y_0$ eine Kurve vom Geschlecht~$2$ und
$\chi\colon Y_1\to Y_0$ eine nichttriviale étale Doppelüberlagerung mit
Deckinvolution $\nu$; die Kurve $Y_1$ hat das Geschlecht~$3$ und ist
hyperelliptisch~\cite[Prop.~4\,i)]{So26}. Bezeichnet $\iota$ ihre
hyperelliptische Involution und $H$ die hyperelliptische Klasse, so gilt
$\omega_{Y_1}\cong\OO(2H)\cong\OO(4W)$ für einen Weierstraß-Punkt $W$. Schränkt
man die Relation $\OO(\Gamma-2\Delta_1)=(1\times\psi_1)^*\mathcal{P}+\pi_2^*L$
aus~\cite[\S5.1]{So26} auf $Y_1\times\{x\}$ ein, so erkennt man, dass
$\psi_1\colon Y_1\to J_1$ durch $x\mapsto[\nu x-x]$ gegeben ist. Die Kurve $Y$
ist eine Zusammenhangskomponente des Faserprodukts $Z$ von $\psi_1$ mit der
Multiplikation mit $2$ auf $J_1$, und $\psi\colon Y\to Y_1$ ist endlich étale vom
Grad~$k$.

Auf $S=Y_1\times Y$ mit den Projektionen $\mathrm{pr}_1$ und $p\colon S\to Y$
seien $D_1$ und $D_2$ die Graphen von $\psi$ bzw.\ $\nu\circ\psi$; sie sind
disjunkt, da $\nu$ keinen Fixpunkt besitzt. Es sei $\cL\in\Pic(S)$ mit
$2\cL=\OO_S(D_1+D_2)$~\cite{MD85}, $\rho\colon X\to S$ die längs $D_1+D_2$
verzweigte Doppelüberlagerung und $\phi=p\circ\rho$; dieser Morphismus ist glatt
mit Fasern vom Geschlecht~$6$. Ferner sei $E_i=\rho^{-1}(D_i)_{\mathrm{red}}$,
sodass $\rho^*D_i=2E_i$ gilt, und $B_i=\OO_X(E_i)$. Für $y\in Y$ setzen wir
$P_1=\psi(y)$, $P_2=\nu P_1$ sowie $\lambda_y=\cL|_{Y_1\times\{y\}}$, eine
Quadratwurzel von $\OO(P_1+P_2)$. Der Punkt $P_1$ heißt \emph{exzeptionell}, wenn
$h^0(P_1+P_2)=2$ gilt; äquivalent dazu ist
$P_2=\iota P_1$~\cite[Prop.~5\,iii)]{So26}. Es gibt vier exzeptionelle Punkte,
und $\Psi=\psi^{-1}(\text{exzeptionelle Punkte})$ besitzt $4k$ Elemente.

Sei $G=\phi_1^*\OO(W)$ mit $\phi_1=\mathrm{pr}_1\circ\rho$. Das von
$G,B_1,B_2,F$ aufgespannte Gitter besitzt die in Tabelle~\ref{tab:lattice}
angegebenen Schnittzahlen; dies sind die Werte aus~\cite[Prop.~8]{So26}, ergänzt
um $B_i\cdot F=1$ ($E_i$ ist ein Schnitt) und $G\cdot F=2$ ($\phi_1$ hat auf den
Fasern den Grad zwei). Die Formel von Riemann--Hurwitz liefert
$\omega=4G+B_1+B_2$ und $\omega^2=12k$. Soul\'es Hilfsbündel lautet
\[
L=\omega-2B_1-2B_2=4G-B_1-B_2,\qquad m=6,\qquad L^2=-20k,
\]
und für allgemeines $y$ gilt
$h^0(L_y)=2$~\cite[Prop.~9\,ii), iii), v)]{So26}.

\begin{table}[ht]
\centering
\begin{tabular}{@{}lcccc@{}}
\toprule
 & $G$ & $B_1$ & $B_2$ & $F$\\
\midrule
$G$ & $0$ & $k$ & $k$ & $2$\\
$B_1$ & $k$ & $-2k$ & $0$ & $1$\\
$B_2$ & $k$ & $0$ & $-2k$ & $1$\\
$F$ & $2$ & $1$ & $1$ & $0$\\
\bottomrule
\end{tabular}
\caption{Schnittzahlen auf der Kodaira-Faserung aus~\cite{So26}.}
\label{tab:lattice}
\end{table}

\subsection{Die Zählung}

\begin{lemma}[$k=4$]\label{lem:k4}
Sei $G_2\subseteq J_1[2]$ die Monodromiegruppe des $J_1[2]$-Torsors $Z\to Y_1$.
Dann gilt $G_2=(\nu_*-1)\,J_1[2]$ und $|G_2|=4$. Folglich ist $k=4$ für jede
Wahl von $Y_0$, von $\chi$ und der Komponente $Y$.
\end{lemma}

\begin{proof}
Der Torsor entsteht durch Rückzug von $[2]\colon J_1\to J_1$ längs $\psi_1$;
seine Monodromie ist daher durch
$\pi_1(Y_1)\to H_1(Y_1,\mathbb{Z})\xrightarrow{(\psi_1)_*}
H_1(J_1,\mathbb{Z})\to J_1[2]$ mit $(\psi_1)_*=\nu_*-1$ gegeben, und jede
Zusammenhangskomponente hat über $Y_1$ den Grad $|G_2|$. Die Ordnung lässt sich
über $\mathbb{C}$ bestimmen. Da die Abbildungsklassengruppe transitiv auf den von
null verschiedenen Klassen in $H^1(Y_0,\mathbb{Z}/2)$ operiert, können wir
$\pi_1(Y_0)=\langle a_1,b_1,a_2,b_2\mid[a_1,b_1][a_2,b_2]\rangle$ wählen und
annehmen, dass $\chi$ durch den Charakter mit $a_1\mapsto 1$ und
$b_1,a_2,b_2\mapsto 0$ gegeben ist. Mit dem Schreier-Repräsentantensystem
$\{1,a_1\}$ wird die Untergruppe $\pi_1(Y_1)$ von $s=a_1^2$, $b_1$, $a_2$, $b_2$
sowie ihren Konjugierten $b_1'$, $a_2'$, $b_2'$ unter $a_1$ erzeugt. Die
definierende Relation und ihre Konjugierte liefern in der Abelschenisierung beide
$b_1'=b_1$, sodass $H_1(Y_1,\mathbb{Z})$ frei mit der Basis
$s,b_1,a_2,b_2,a_2',b_2'$ ist. Die Involution $\nu$ entspricht der Konjugation
mit $a_1$; auf der Homologie lässt sie $s$ und $b_1$ fest und vertauscht
$a_2\leftrightarrow a_2'$ sowie $b_2\leftrightarrow b_2'$. Das Bild von $\nu_*-1$
wird von $a_2'-a_2$ und $b_2'-b_2$ aufgespannt, die sich zu einer Basis ergänzen
lassen; seine Reduktion modulo~$2$ besitzt daher die Ordnung~$4$.
\end{proof}

\begin{lemma}[effektive Quadratwurzeln]\label{lem:keff}
Sei $P_1$ exzeptionell. Die Klassen $\lambda_y$, $y\in\psi^{-1}(P_1)$, bilden
eine Nebenklasse von $G_2$ in der Menge der Quadratwurzeln von $H$. Jede
Nebenklasse von $G_2$ in dieser Menge enthält höchstens zwei effektive Klassen.
\end{lemma}

\begin{proof}
Für $(x,z)\in Z$ ist die Klasse $x+z\in\Pic^1(Y_1)$ wegen $2(x+z)=2x+\nu x-x$
eine Quadratwurzel von $x+\nu x$. Die Abbildung
$y\mapsto\lambda_y-(\psi(y)+\psi_2(y))$ mit der zweiten Projektion
$\psi_2\colon Y\to J_1$ nimmt Werte in $J_1[2]$ an und ist daher auf der
zusammenhängenden Kurve $Y$ konstant. Über $P_1$ bilden die Punkte von $Y$ eine
$G_2$-Bahn in der Faser von $Z$; folglich bilden die $\lambda_y$ eine Nebenklasse
von $G_2$. Eine effektive Quadratwurzel von $H$ hat die Gestalt $\OO(Q)$ mit
$2Q\in|H|$, gehört also zu einem Weierstraß-Punkt. Seien $W\ne W'$ effektiv in
derselben Nebenklasse. Dann gilt $W'-W=(\nu-1)g$ für ein $g\in J_1[2]$, also
$(1+\nu)(W'-W)=(\nu^2-1)g=0$ und damit $W'+\nu W'\sim W+\nu W$. Wären diese
Divisoren verschieden, so würde sich $W+\nu W$ in einem Büschel bewegen und läge
folglich in $|H|$; daraus ergäbe sich $\nu W=\iota W=W$, im Widerspruch zur
Fixpunktfreiheit von $\nu$. Also gilt $W'=\nu W$, und eine Nebenklasse enthält
höchstens die beiden effektiven Wurzeln $W,\nu W$.
\end{proof}

\begin{lemma}[Torsion der Länge eins]\label{lem:length}
Der Träger der Torsion von $R^1p_*\cL$ besteht genau aus den Punkten $y$ mit
$h^0(\lambda_y)\ne 0$, und in jedem dieser Punkte hat die Torsion die Länge eins.
\end{lemma}

\begin{proof}
Generisch gilt $h^0(\lambda_y)=0$~\cite[Prop.~5\,ii)]{So26}, also $p_*\cL=0$.
Es sei $\pi$ eine Uniformisierende in $y$. Die exakte Sequenz
$0\to\cL\xrightarrow{\pi}\cL\to\lambda_y\to 0$ liefert
$R^1p_*\cL[\pi]\cong H^0(\lambda_y)$; die Torsion $T_y$ ist somit genau dann von
null verschieden, wenn $h^0(\lambda_y)\ne 0$ gilt, und wegen
$h^0(\lambda_y)\le 1$ ist sie zyklisch. Ebenso ist $R^1p_*\cL[\pi^2]$ isomorph
zum Raum der Schnitte von $\cL$ über $Y_1\times\operatorname{Spec}\OO_{Y,y}/\pi^2$;
daher gilt $\ell(T_y)\ge 2$ genau dann, wenn sich der Schnitt $s_W$ von
$\lambda_y=\OO(W)$ in erster Ordnung fortsetzen lässt. Dies ist genau dann der
Fall, wenn $v\cup s_W=0$ in $H^1(\OO(W))$ gilt, wobei $v\in H^1(\OO_{Y_1})$ den
Tangentialvektor der Abbildung $y\mapsto\lambda_y$ bezeichnet, wenn also $v$ auf
der Geraden $d\mathrm{AJ}_W(T_WY_1)$ liegt. Unter der Dualität $H^1(\OO_{Y_1})\cong
H^0(\omega_{Y_1})^\vee$ wird diese Gerade von der Auswertung in $W$ aufgespannt.
Da $2\lambda_y=\OO(\psi(y)+\nu\psi(y))$ gilt und $\psi$ étale ist, ist $2v$ das
Funktional $\alpha\mapsto\alpha_{P_1}(t)+\alpha_{P_2}(d\nu\,t)$ für einen von
null verschiedenen Tangentialvektor $t$ in $P_1$. Für einen exzeptionellen Punkt
gilt $P_2=\iota P_1$; da $\iota$ mit $\nu$ kommutiert, hat die Involution
$\nu\iota$ den Fixpunkt $P_2$ mit Ableitung $-1$, sodass $d\nu=-d\iota$ in $P_1$
folgt. Wegen $\iota^*\alpha=-\alpha$ für alle $\alpha\in H^0(\omega_{Y_1})$
ergibt sich $2v=2\,\mathrm{ev}_{P_1,t}\ne 0$. Die Auswertungen in $P_1$ und $W$
sind nur dann proportional, wenn die kanonische Abbildung $P_1$ und $W$
identifiziert, wenn also $P_1\in\{W,\iota W\}=\{W\}$ gilt. Aus $P_1=W$ folgte
jedoch $W\sim P_2$ und damit $P_2=P_1$, ein Widerspruch. Somit ist $\ell(T_y)=1$.
\end{proof}

\begin{proposition}[Gradformel]\label{prop:degree}
Es sei $T$ die Anzahl der Punkte $y\in Y$ mit $h^0(\lambda_y)\ne 0$; sie liegen
sämtlich in $\Psi$. Dann gilt
\[
\deg\phi_*L=-7k+T .
\]
\end{proposition}

\begin{proof}
Es sei $K_1=\mathrm{pr}_1^*\omega_{Y_1}=\omega_{S/Y}$. Die Formel für das
kanonische Bündel einer Doppelüberlagerung~\cite[I.17]{BHPV04} liefert
$\omega_{X/Y}=\rho^*(K_1+\cL)$, und es gilt
$B=2B_1+2B_2=\rho^*(D_1+D_2)=\rho^*(2\cL)$. Also ist $L=\rho^*(K_1-\cL)$, und
mit $\rho_*\OO_X=\OO_S\oplus\cL^{-1}$~\cite[(13)]{So26} folgt
\[
\phi_*L=p_*(K_1-\cL)\oplus p_*(K_1-2\cL).
\]
\emph{Zweiter Summand.} Da $x+\iota x\in|H|$ für jedes $x$ gilt, liefert das
Seesaw-Prinzip $\OO_S(D_i+\iota D_i)\cong\mathrm{pr}_1^*\OO(H)\otimes p^*M_i$ mit
$M_i\in\Pic(Y)$, wobei $\iota D_i=(\iota\times 1)(D_i)$ gesetzt ist. Schneiden
mit $\{x\}\times Y$, das jeden der Divisoren $D_i$ und $\iota D_i$ transversal in
$k$ Punkten trifft und auf dem $\mathrm{pr}_1^*\OO(H)$ trivial ist, ergibt
$\deg M_i=2k$. Wegen $K_1=\mathrm{pr}_1^*\OO(2H)$ gilt
\[
K_1-2\cL=K_1-D_1-D_2\cong\OO_S(\iota D_1+\iota D_2)\otimes p^*(-M_1-M_2).
\]
Der Automorphismus $\iota\times 1$ von $S$ über $Y$ und~\cite[(20)]{So26} liefern
$p_*\OO_S(\iota D_1+\iota D_2)\cong p_*\OO_S(D_1+D_2)\cong\OO_Y$, also
$\deg p_*(K_1-2\cL)=-4k$.

\emph{Erster Summand.} Relative Dualität ergibt
$p_*(K_1-\cL)\cong\mathcal{H}om(R^1p_*\cL,\OO_Y)$, nach Lemma~\ref{lem:length}
also $\deg p_*(K_1-\cL)=-(\deg R^1p_*\cL-T)$. Wegen $p_*\cL=0$ liefert der Satz
von Grothendieck--Riemann--Roch für $p$ die Gleichung
$\deg R^1p_*\cL=-\bigl(\tfrac12\cL^2-\tfrac12K_1\cdot\cL\bigr)$, und aus
$4\cL^2=(D_1+D_2)^2=-8k$ sowie $2K_1\cdot\cL=K_1\cdot(D_1+D_2)=8k$ folgt
$\deg R^1p_*\cL=3k$. Damit ist $\deg p_*(K_1-\cL)=T-3k$.

Addition der beiden Grade ergibt $\deg\phi_*L=T-3k-4k$.
Nach~\cite[Prop.~5\,ii)]{So26} erzwingt $h^0(\lambda_y)\ne 0$, dass $P_1$
exzeptionell ist; die in $T$ gezählten Punkte liegen daher in $\Psi$.
\end{proof}

\begin{theorem}\label{thm:soule-count}
Für die Kodaira-Faserung aus~\cite{So26} mit $L=\omega-2B_1-2B_2$ gilt
\[
\deg\phi_*L=-7k+4k_{\mathrm{eff}},\qquad k=4,\qquad 0\le k_{\mathrm{eff}}\le 2,
\]
wobei $4k_{\mathrm{eff}}=T$ gesetzt ist. Insbesondere gilt
$\deg\phi_*L\le-20<-12=-3k$; somit ist Proposition~9\,vi) in arXiv:2611.05217v1
falsch. Ferner ist $q=-\deg\phi_*L/\omega^2\ge 5/12$, und der führende
Koeffizient $5/3-6q$ in~\cite[(36)]{So26} erfüllt $5/3-6q\le-5/6$.
\end{theorem}

\begin{proof}
Proposition~\ref{prop:degree} liefert $\deg\phi_*L=-7k+T$. Nach
Lemma~\ref{lem:k4} ist $k=4$, und Lemma~\ref{lem:keff} zeigt, dass jede der vier
Fasern von $\psi$ über einem exzeptionellen Punkt höchstens zwei Punkte zu $T$
beiträgt, also $T\le 8$. Damit folgt $\deg\phi_*L\le-28+8=-20$ sowie
$q\ge 20/48=5/12$.
\end{proof}

\begin{theorem}[strukturelle Fassung]\label{thm:soule-structural}
Unter den Voraussetzungen von~\cite{So26} ist das Linearsystem $|L_\eta|$ für
jede Wahl der Komponente $Y$ und von $\cL$ basispunktfrei. Folglich gilt
$\deg\phi_*L\le L^2/6=-10k/3<-3k$ und $M(L)\le 0$; insbesondere erfüllt keine
Wahl der Komponente $Y$ oder von $\cL$ die Aussage von Proposition~9\,vi).
\end{theorem}

\begin{proof}
\emph{Schritt 1: Fasern außerhalb von $\Psi$.} Sei $y\notin\Psi$, also
$P_2\ne\iota P_1$. Nach~\cite[Prop.~5\,ii)]{So26} ist dann $h^0(\lambda_y)=0$,
und der Satz von Riemann--Roch liefert $h^0(\omega_{Y_1}-\lambda_y)=1$ sowie
$h^0(\omega_{Y_1}-2\lambda_y)=h^0(P_1+P_2)=1$. Wegen $\omega_{Y_1}\cong\OO(2H)$
und $P_i+\iota P_i\in|H|$ ist $\iota P_1+\iota P_2$ der einzige effektive Divisor
in $|\omega_{Y_1}-2\lambda_y|$. Nach der im Beweis von
Proposition~\ref{prop:degree} hergeleiteten Zerlegung von $\phi_*L$ wird
$H^0(X_y,L_y)$ von $\rho^*\sigma_y$ und $t_y\,\rho^*\tau_y$ aufgespannt; dabei
erzeugt $\sigma_y$ den Raum $H^0(\omega_{Y_1}-\lambda_y)$, $\tau_y$ besitzt den
Divisor $\iota P_1+\iota P_2$, und $t_y$ ist der in den beiden
Verzweigungspunkten verschwindende Schnitt von $\rho^*\lambda_y$. Folglich hat
$|L_y|$ genau dann einen Basispunkt, wenn $\sigma_y$ in einem der Punkte $P_1$,
$P_2$, $\iota P_1$, $\iota P_2$ verschwindet, und nach Riemann--Roch verschwindet
$\sigma_y$ in einem Punkt $Q$ genau dann, wenn $h^0(\lambda_y+Q)\ne 0$ gilt.

\emph{Schritt 2: Spezialisierung.} Auf der offenen Teilmenge $Y^\circ\subseteq Y$
der Punkte mit $h^0(\lambda_y)=0$ gilt $h^0(\omega_{Y_1}-\lambda_y)=1$; dort ist
$p_*(K_1-\cL)$ also ein Geradenbündel, dessen Bildung mit Basiswechsel
vertauscht. Für jede der vier Abbildungen $Q=\psi,\nu\psi,\iota\psi,\iota\nu\psi$
von $Y$ nach $Y_1$ ist daher die Menge der $y\in Y^\circ$ mit
$\sigma_y(Q(y))=0$ abgeschlossen in $Y^\circ$. Keine dieser Mengen stimmt mit
ganz $Y^\circ$ überein: Sei nämlich $x_0$ exzeptionell; nach den
Lemmata~\ref{lem:k4} und~\ref{lem:keff} existiert ein $y_0\in\psi^{-1}(x_0)$ mit
$h^0(\lambda_0)=0$, wobei $\lambda_0=\lambda_{y_0}$ gesetzt sei. Wegen $\nu x_0=\iota x_0\ne
x_0$ liegen die vier Punkte $Q(y_0)$ in $\{x_0,\iota x_0\}$, keiner von ihnen ist
ein Weierstraß-Punkt, und es gilt $2\lambda_0\sim x_0+\iota x_0\sim H$. Gilt
$\lambda_0+q\sim a+b$ mit $q\in\{x_0,\iota x_0\}$, so folgt $2a+2b\in|H+2q|$. Da
$q$ kein Weierstraß-Punkt ist, gilt $h^0(H+2q)=2$ und $|H+2q|=|H|+2q$, also
$2a+2b=c+\iota c+2q$ für einen Punkt $c$. Auf der linken Seite sind sämtliche Vielfachheiten gerade,
während für $c\ne\iota c$ auf der rechten Seite der Punkt $c$ oder der Punkt $q$
mit ungerader Vielfachheit auftritt. Also folgt $c=\iota c$ und $a+b=c+q$, womit
$\lambda_0\sim c$ effektiv wäre, ein Widerspruch. Daher verschwindet
$\sigma_{y_0}$ in keinem der Punkte $Q(y_0)$; die vier abgeschlossenen Teilmengen
sind somit endlich, und nach Schritt~1 ist $|L_y|$ für alle bis auf endlich viele
$y\in Y\setminus\Psi$ basispunktfrei.

\emph{Schritt 3: Die generische Faser.} Für $y\in Y\setminus\Psi$ gilt
$h^0(L_y)=2$, sodass $\phi_*L\otimes k(y)\to H^0(X_y,L_y)$ bijektiv ist. Der
Abschluss eines Basispunkts von $|L_\eta|$ bildet eine horizontale Kurve, längs
der $\phi^*\phi_*L\to L$ nicht surjektiv ist; sie schneidet jede Faser $X_y$, und
für $y\in Y\setminus\Psi$ trifft sie $X_y$ in einem Basispunkt von $|L_y|$, im
Widerspruch zu Schritt~2. Also ist $|L_\eta|$ basispunktfrei, und
Satz~\ref{thm:rank2} liefert mit $m=6$ und $L^2=-20k$ die behauptete Schranke.

\emph{Schritt 4: Bezug zu~\cite{So26}.} In~\cite[\S6.3.3]{So26} wird die
Quadratwurzel $\cL'$ in Abhängigkeit vom Punkt $\bar y'\notin\Phi'$ gewählt,
und~\cite[Lemma~3]{So26} zeigt, dass $L'_{P'}$ keinen Basispunkt besitzt. Da
$h^0(L'_{P'})=2$ mit dem Rang von $\phi'_*L'$ übereinstimmt, erzwingt dies nach
dem Argument aus Schritt~3 bereits die Basispunktfreiheit von $|L'_\eta|$; jedes
Bündel, auf das~\cite[(28)]{So26} angewandt wird, fällt somit unter
Satz~\ref{thm:rank2}. Die Schritte~1--3 übertragen sich wörtlich nach dem
Basiswechsel $a\colon Y'\to Y$ aus~\cite[\S6.3.1]{So26}, wenn man $\psi$ durch
$\psi\circ a$ ersetzt. Es sei $J'$ die Jacobische von $Y'$. Da
$\operatorname{Hom}(J_1,J')$ torsionsfrei ist, gilt
$\Pic(Y_1\times Y')[2]=\mathrm{pr}_1^*J_1[2]\oplus\mathrm{pr}_2^*J'[2]$; auf den
Fasern schränkt sich $\cL'$ daher auf die Klassen $\lambda_{a(y')}+\varepsilon$ für
ein festes $\varepsilon\in J_1[2]$ ein. Über einem exzeptionellen Punkt $x_0$ bilden
die Klassen $\lambda_y+\varepsilon$, $y\in\psi^{-1}(x_0)$, eine Nebenklasse von
$G_2$; nach Lemma~\ref{lem:keff} gibt es folglich ein $y_0\in\psi^{-1}(x_0)$ mit
$h^0(\lambda_{y_0}+\varepsilon)=0$, und wegen der Surjektivität von $a$ lässt sich
$y_0$ zu einem Punkt in $(\psi\circ a)^{-1}(x_0)$ liften; mehr wird in Schritt~2
nicht benötigt. Lemma~\ref{lem:k4} geht nur für $\psi$ mit $k=4$ ein. Die
Basispunktfreiheit ist daher unabhängig von der punktweisen Wahl.
\end{proof}

\begin{remark}[Stelle, an der der Beweis von Proposition~9 scheitert]\label{rem:gap}
Ungleichung~(24) aus~\cite{So26} besagt
$\deg(R^1\phi_*B)_{\mathrm{tor}}\ge 8k$. Sie stützt sich
auf~\cite[Prop.~9\,iv)]{So26}, wonach $h^0(B_{\bar y})=3$ gilt, sobald $P_1$
exzeptionell ist. Der dortige Beweis verwendet den Satz „when $P_1$ is
exceptional $h^0(\lambda)=1$“ unter Berufung auf die Teile ii) und iii)
von~\cite[Prop.~5]{So26}. Proposition~5\,ii) besagt jedoch, dass $h^0(\lambda)=0$
gilt, sofern $\lambda$ kein Weierstraß-Punkt ist, und dass in diesem Ausnahmefall
$h^0(\lambda)=1$ sowie $h^0(P_1+P_2)=2$ folgen. Bewiesen wird dort also die
Implikation „$\lambda$ effektiv $\Rightarrow$ $P_1$ exzeptionell“, d.\,h.\ die
Umkehrung der benötigten Aussage. Nach Lemma~\ref{lem:keff} versagt die benutzte
Implikation in mindestens $4k-8=8$ der $16$ Punkte von $\Psi$; dort gilt
$h^0(B_{\bar y})=2$ und $h^0(L_{\bar y})=3$, sodass Proposition~9\,iv) und~v) in
diesen Punkten nicht zutreffen. Zusammen mit~\cite[(16), (18), (22)]{So26}
liefert Proposition~\ref{prop:degree} daher
$\deg(R^1\phi_*B)_{\mathrm{tor}}=4k+T\le 24$ anstelle von $\ge 32$. Da eine
globale Abänderung von $\cL$ um ein Element aus $J_1[2]$ sämtliche Nebenklassen
gleichzeitig verschiebt, ändert keine Wahl von $\cL$ die Schranke aus
Lemma~\ref{lem:keff}. Theorem~1 aus~\cite{So26} bleibt als Aussage offen;
hinfällig ist der Beweis über Proposition~9\,vi), und nach
Korollar~\ref{cor:rank2} in Verbindung mit Schritt~4 des Beweises von
Satz~\ref{thm:soule-structural} lässt sich der Büschelschritt auch nicht durch
ein anderes Hilfsbündel $L'$ vom Rang zwei mit basispunktfreiem $|L'_\eta|$
retten. Besitzt $|L'_\eta|$ hingegen einen Basispunkt, so ist
$\phi_*L'\otimes k(P)\to H^0(X_P,L'_P)$ für alle $P$ außerhalb einer endlichen
Ausnahmemenge bijektiv; nach dem Argument aus Schritt~3 jenes Beweises haben dann
für alle derartigen $P$ je zwei Schnitte von $L'_P$ eine gemeinsame Nullstelle,
sodass die Büschelvoraussetzung in~\cite[(28)]{So26} unerfüllbar ist.
\end{remark}

\section{Reichweite: Kreuzterme, einzelne Schnitte, Quellbündel}\label{sec:scope}

Dieser Abschnitt stellt drei weitere Folgerungen aus dem Nef-Lemma und aus der
Identität hinter Satz~\ref{thm:rankr} zusammen. Die formalen Aussagen sind
geometrischer Natur und vollständig bewiesen; die Bemerkungen zu ihrer
arithmetischen Bedeutung dienen der Einordnung.

\subsection{Kreuzterme}
Proposition~3 aus~\cite{So26} ist für zwei Geradenbündel $L_1,L_2$ mit Fasergraden
$d_1,d_2$ formuliert: $-\Lb_1\cdot\Lb_2\le d_2h(s_1)+d_1h(s_2)$. Stammt $s_1$ aus
der Harder--Narasimhan-Stufe $a$ von $\phi_*L_1$ und $s_2$ aus der Stufe $b$ von
$\phi_*L_2$, wobei die Höhen durch die Steigungen $\mu_{1,a}$ und $\mu_{2,b}$
bestimmt sind, so lautet der führende Koeffizient
\[
M(L_1,L_2)=\frac{-L_1\cdot L_2+d_2\mu_{1,a}+d_1\mu_{2,b}}{\omega^2}
=-\frac{(L_1-\mu_{1,a}F)\cdot(L_2-\mu_{2,b}F)}{\omega^2}.
\]

\begin{proposition}\label{prop:cross}
Besitzen die horizontalen Basisdivisoren $B_{1,a}$ und $B_{2,b}$ keine gemeinsame
Komponente, so gilt $(L_1-\mu_{1,a}F)\cdot(L_2-\mu_{2,b}F)\ge 0$ und folglich
$M(L_1,L_2)\le 0$.
\end{proposition}

\begin{proof}
Man schreibe $L_1-\mu_{1,a}F=N_{1,a}+B_{1,a}+V_{1,a}$ und verfahre analog für
$L_2$, wobei $N_{1,a},N_{2,b}$ nach Lemma~\ref{lem:nef} nef sind. Beim
Ausmultiplizieren ist jeder Summand nichtnegativ: das Schnittprodukt zweier nef
Divisoren, das Produkt eines nef mit einem effektiven Divisor,
$B_{1,a}\cdot B_{2,b}\ge 0$ (da keine gemeinsame Komponente vorliegt), das
Produkt eines horizontalen mit einem vertikalen Divisor sowie
$V_{1,a}\cdot V_{2,b}=0$.
\end{proof}

\subsection{Einzelne Schnitte und Jets längs Kurven}
Für einen einzelnen Schnitt $s$ und eine irreduzible horizontale Kurve $C$, die
nicht im Divisor von $s$ enthalten ist, liefert der Vergleich punktweiser Normen
mit den Supremumsnormen eine Abschätzung der Gestalt $-\Lb\cdot\overline C\le(C\cdot
F)\,h(s)$, deren geometrischer Koeffizient für $s$ aus $E_j$ gleich
$-(L-\mu_jF)\cdot C/\omega^2$ ist. Normale Jets der Ordnung $b$ längs einer Kurve
$C$, die in $B_j$ mit der Vielfachheit $b$ auftritt, sind Schnitte von
$(L-bC)|_C$.

\begin{proposition}\label{prop:curves}
Es sei $C$ eine irreduzible horizontale Kurve. (a) Ist $C\not\subseteq B_j$, so
gilt $(L-\mu_jF)\cdot C\ge 0$. (b) Tritt $C$ in $B_j$ mit der Vielfachheit $b$
auf, so gilt $(L-bC-\mu_jF)\cdot C\ge 0$.
\end{proposition}

\begin{proof}
(a) Es gilt $L-\mu_jF=N_j+B_j+V_j$ mit $N_j\cdot C\ge 0$, $B_j\cdot C\ge 0$ und
$V_j\cdot C\ge 0$. (b) Es ist $L-bC-\mu_jF=N_j+(B_j-bC)+V_j$, wobei $C$ keine
Komponente von $B_j-bC$ bildet.
\end{proof}

\subsection{Quellbündel, die nicht von der Basis zurückgezogen sind}
Die Identität aus Satz~\ref{thm:rankr} besitzt eine Verallgemeinerung für ein
beliebiges Quellbündel.

\begin{proposition}\label{prop:universal}
Es sei $\cF$ lokal frei vom Rang $r\ge 2$ auf $X$ und $\cF\to L$ ein Morphismus
mit Bild $A\otimes\cI_Z$, wobei $A$ ein Geradenbündel ist und $Z$ endliche Länge
besitzt. Dann ist der Kern $K$ lokal frei vom Rang $r-1$, und es gilt
\[
\Bigl(A-\frac{c_1(\cF)}{r}\Bigr)^2=\frac{\Delta(K)+2(r-1)\ell(Z)}{r}-\frac{r-1}{r^2}\,\Delta(\cF).
\]
\end{proposition}

\begin{proof}
Es gilt $c_1(K)=c_1(\cF)-A$ sowie $\ch_2(K)=\ch_2(\cF)-\frac12A^2+\ell(Z)$. Setzt
man dies in $\Delta(K)=c_1(K)^2-2(r-1)\ch_2(K)$ ein und verwendet
$\ch_2(\cF)=(c_1(\cF)^2-\Delta(\cF))/(2r)$, so folgt durch Zusammenfassen die
Behauptung.
\end{proof}

Für $\cF=\phi^*E$ gilt $\Delta(\cF)=0$ und $c_1(\cF)\equiv(\deg E)F$, womit sich
die Proposition auf Satz~\ref{thm:rankr} reduziert. Ist $\Delta(K)\ge 0$, so
kann allein ein positiver Beitrag $\Delta(\cF)>0$ die linke Seite negativ machen;
hierfür ist ein Bündel erforderlich, das nicht von der Basis zurückgezogen ist.
Im arithmetischen Rahmen führt dies zurück auf Parshins Programm einer
arithmetischen Bogomolov--Miyaoka--Yau-Ungleichung~\cite{Pa88,MB90} sowie auf
arithmetische Ungleichungen vom Bogomolov--Gieseker-Typ~\cite{Mo94}.
Soul\'e~\cite{So05} verfolgte Miyaokas Zugang über Erweiterungen vom Rang zwei;
er hält dort fest, dass die arithmetische zweite Chern-Zahl „contains an
archimedean summand -- an integral over the set of complex points of $X$ -- which
is not easy to bound from above“, und gewinnt untere Schranken für einzelne
sukzessive Minima anstelle einer oberen Schranke für $\ob^{\,2}$. Weitere
Schlüsse ziehen wir daraus nicht.

\section{Schlussbemerkungen}\label{sec:conclusion}

Innerhalb der hier betrachteten Klasse von Argumenten beruht jede verfügbare
Positivität auf demselben Prinzip: Das Schnittprodukt eines nef Divisors mit
einem effektiven Divisor ist nichtnegativ. In jedem oben behandelten Fall geht
diese Positivität mit genau jenem Vorzeichen in den führenden Koeffizienten ein,
das den Spielraum nichtpositiv macht, während ein PMB-Argument einen positiven
Spielraum erfordert. Für Büschel zeigen dies Satz~\ref{thm:rank2} in Rang zwei
und Korollar~\ref{cor:ss} für semistabiles $\phi_*L$ in jedem Rang. Offen bleibt
in dieser Richtung die Frage einer arithmetischen Anomalie für instabiles
$\phi_*L$ vom Rang $r\ge 3$ (Offenes Problem~\ref{prob:anomaly}). Für~\cite{So26},
Version~1, ist die Folgerung eng umgrenzt und konkret: Proposition~9\,vi) ist
falsch, Ungleichung~(36) besitzt einen negativen führenden Koeffizienten, und der
dort gegebene Beweis von Theorem~1 ist unvollständig.

\phantomsection
\addcontentsline{toc}{section}{KI-Offenlegung}
\begingroup
\sloppy\emergencystretch=2em
% =============================================================================
% AI DISCLOSURE -- PIPELINE-STANDARD, DEUTSCHE FASSUNG (seit 2026-08-17, T-20260817-02)
% =============================================================================
% Sprachreine Fassung fuer rein deutschsprachige Manuskripte. Text 1:1 aus dem
% deutschen Block von ai_disclosure_STANDARD.tex uebernommen (keine inhaltliche
% Aenderung) -- nur die Ueberschrift ist jetzt sprachrein und der englische
% Block entfaellt. Die bilinguale Datei ai_disclosure_STANDARD.tex bleibt
% unveraendert bestehen (Bestandsverweise duerfen weiter darauf zeigen).
%
% ABWEICHEN NACH UNTEN: Wurde in einem Paper tatsaechlich deutlich weniger
% KI eingesetzt, wird NICHT auf eine andere Datei gewechselt, sondern die
% Formulierung hier abgeschwaecht -- z. B. "in erheblichem Masse in allen
% Phasen" ersetzen durch die tatsaechlich zutreffenden Phasen, und aus der
% Qualitaetssicherungs-Aufzaehlung streichen, was nicht stattfand.
% Der Grundsatz bleibt: ehrlich benennen, was war -- nicht abgrenzen.
%
% Verwendung (nur fuer deutschsprachige Manuskripte):
%   \input{../../_templates/disclosure/ai_disclosure_STANDARD_de}
% =============================================================================

\section*{KI-Offenlegung}

\noindent
KI-Systeme wurden in erheblichem Ma\ss{}e in allen Phasen dieser Arbeit
eingesetzt: Konzeption und Forschungsdiskussion, Literatur- und
Quellenarbeit, mathematische Analyse und Beweisarbeit, Erstellung und
Pr\"ufung der Rechenskripte, Ausf\"uhrung und Auswertung der
Berechnungsl\"aufe, Text, \"Ubersetzung und Satz.

\medskip

\noindent
Qualit\"atssicherung: Einsatz mehrerer unabh\"angiger Modelle
verschiedener Anbieter mit Gegenpr\"ufungen und adversarialen KI-Reviews;
nachvollziehbare, versionierte Rechenskripte mit gespeicherten
Ergebnisdateien und Pr\"ufsummen; fortlaufende Beweis- und
Registerdokumente; strikte Claim-Hygiene (Behauptungen nur mit
dokumentiertem Beleg, Nachtr\"age statt stiller Korrekturen); Changelogs
\"uber alle Versionen.

\endgroup

\phantomsection
\addcontentsline{toc}{section}{Literatur}
\begin{thebibliography}{99}
\small
\frenchspacing

\bibitem{BHPV04}
W.~P. Barth, K.~Hulek, C.~A.~M. Peters, A.~Van de Ven,
\emph{Compact Complex Surfaces}, 2nd enlarged ed.,
Ergebnisse der Mathematik und ihrer Grenzgebiete (3) 4, Springer, Berlin, 2004.

\bibitem{Bo78}
F.~A. Bogomolov,
Holomorphic tensors and vector bundles on projective varieties,
\emph{Izv. Akad. Nauk SSSR Ser. Mat.} 42 (1978), 1227--1287;
engl.\ Übers.\ \emph{Math. USSR-Izv.} 13 (1979), 499--555.

\bibitem{CH88}
M.~Cornalba, J.~Harris,
Divisor classes associated to families of stable varieties, with applications
to the moduli space of curves,
\emph{Ann. Sci. \'Ec. Norm. Sup\'er.} (4) 21 (1988), 455--475.

\bibitem{HN75}
G.~Harder, M.~S. Narasimhan,
On the cohomology groups of moduli spaces of vector bundles on curves,
\emph{Math. Ann.} 212 (1975), 215--248.

\bibitem{HL10}
D.~Huybrechts, M.~Lehn,
\emph{The Geometry of Moduli Spaces of Sheaves}, 2nd ed.,
Cambridge Mathematical Library, Cambridge University Press, Cambridge, 2010.

\bibitem{Ko67}
K.~Kodaira,
A certain type of irregular algebraic surfaces,
\emph{J. Analyse Math.} 19 (1967), 207--215.

\bibitem{MD85}
M.~Martin-Deschamps,
La construction de Kodaira--Parshin,
in: \emph{S\'eminaire sur les pinceaux arithm\'etiques: la conjecture de Mordell},
Ast\'erisque 127 (1985), 261--273.

\bibitem{Mi77}
Y.~Miyaoka,
On the Chern numbers of surfaces of general type,
\emph{Invent. Math.} 42 (1977), 225--237.

\bibitem{MB90}
L.~Moret-Bailly,
Hauteurs et classes de Chern sur les surfaces arithm\'etiques,
in: \emph{S\'eminaire sur les pinceaux de courbes elliptiques (\`a la recherche
de \guillemotleft Mordell effectif\guillemotright)}, Ast\'erisque 183 (1990), 37--58.

\bibitem{Mo94}
A.~Moriwaki,
Inequality of Bogomolov--Gieseker type on arithmetic surfaces,
\emph{Duke Math. J.} 74 (1994), 713--761.

\bibitem{Pa88}
A.~N. Parshin,
The Bogomolov--Miyaoka--Yau inequality for the arithmetical surfaces and its
applications,
in: \emph{S\'eminaire de Th\'eorie des Nombres, Paris 1986--87},
Progr. Math. 75, Birkh\"auser, Boston, 1988, 299--312.

\bibitem{So05}
C.~Soul\'e,
Semi-stable extensions on arithmetic surfaces,
arXiv:math/0505079 (2005); in: \emph{Moduli Spaces and Arithmetic Geometry
(Kyoto, 2004)}, Adv. Stud. Pure Math. 45 (2006), 283--295.

\bibitem{So26}
C.~Soul\'e,
On the $abc$ conjecture,
arXiv:2611.05217, Version~1 (3.~November 2026).

\bibitem{Xi87}
G.~Xiao,
Fibered algebraic surfaces with low slope,
\emph{Math. Ann.} 276 (1987), 449--466.

\end{thebibliography}

\end{document}
